% ID: SEL-D0008-BB-E04
% ESTADO: validada. S009, 2026-11-07; tres pasadas del mismo autor y revisión matemática, didáctica y visual.
\textbf{a) Mínimo y unidades.}\par\nopagebreak[4]
La cantidad producida es no negativa en el contexto del problema. Completando el cuadrado,
\[
f(x)=x^2-6x+10=(x-3)^2+1\geq1.
\]
La igualdad sólo se alcanza en $x=3$. La producción óptima es
\[
\boxed{3\text{ miles de kg}=3000\text{ kg semanales}},
\]
y el coste mínimo es
\[
\boxed{1\text{ millar de euros}=1000\text{ euros}.}
\]

También se comprueba con $f'(x)=2x-6$: negativa antes de 3 y positiva después. El mínimo está dentro del dominio contextual $x\geq0$.

\textbf{b) Tangente en x=4 y representación.}\par\nopagebreak[4]
\[
f(4)=16-24+10=2,\qquad f'(4)=2.
\]
La recta pasa por $(4,2)$ con pendiente 2:
\[
y-2=2(x-4)\quad\Longrightarrow\quad\boxed{y=2x-6}.
\]
Aquí $x$ se mide en miles de kg e $y$ en miles de euros. La parábola abre hacia arriba, tiene vértice $(3,1)$ y corte vertical $(0,10)$; no corta el eje horizontal. La tangente toca la curva en $(4,2)$.
\begin{center}\includegraphics[width=0.62\linewidth]{figuras/s009-coste.pdf}\end{center}
Comprobación algebraica: $f(x)-(2x-6)=(x-4)^2\geq0$, con igualdad sólo en 4. La gráfica muestra $x\geq0$, que es el rango físico de producción; la fórmula polinómica está definida en todo $\mathbb R$.
